\section{Materiales y métodos}
 Se estudió la transición de fase ferromagnética-paramagnética de una ferrita mixta de Mn y Zn, cuyas proporciones son desconocidas, esto con el objetivo de determinar la temperatura crítica de la muestra y así poder caracterizar la composición de la misma.
 El dispositivo experimental (ver figura)consiste en dos bobinas enrolladas alrededor de la muestra de ferrita, ambas con N vueltas y de inductancia L. A la primera bobina se encuentra conectada una fuente de corriente alterna (asumo sinometer pero hoy pregunto) que se utilizó para crear el campo magnetizante \(H\) que luego induciría el campo magnético \(M\) sobre la muestra. Para poder medir una cantidad proporcional a \(H\), se introdujo una resistencia de (resistencia) entre el generador de señales y la bobina, de manera que se midió una tensión \(V_1\) proporcional a la corriente \(I_1\) que genera el campo magnetizante. Por otro lado, el flujo magnético generado por el campo inducido \(M\), induce a su vez una tensión \(V_2\) en la segunda bobina, que es proporcional al cambio en el flujo. Por esto, se utilizó como integrador un circuito RC de resistencia R y capacitancia C, de manera que, bajo la condición \(R \gg X_c\), al medir la tensión \(V_c\) sobre el capacitor, se obtiene una cantidad proporcional a la magnetización \(M\) de la muestra. Finalmente, para obtener un gráfico del ciclo de histéresis, se utilizó una interfaz (uomo pero hoy me fijo bien) que permitió medir simultáneamente las tensiones \(V_1\) y \(V_c\), y así graficar \(M\) en función de \(H\).
 Asimismo, para poder observar la transición de fase, se utilizó un horno (cerámico???) que permitió aumentar la temperatura de la muestra en un rango de (aca el rango). Para poder medir la temperatura de la muestra, se utilizó una termocupla (marca) que se encuentra en contacto directo con la muestra, permitiendo así obtener pares de datos (m y t no se como ponerlo). 
    %aca va la figura jeje
    A partir de los datos obtenidos, se graficó la magnetización \(M\) en función de la temperatura \(T\). Para determinar un valor de la temperatura crítica \(T_c\), se realizó un ajuste de los datos en base a la Ley de Bloch dada por
    \begin{equation}
        M(T) = 1 - \left(\frac{T}{T_c}\right)^{n}
        \label{eq:bloch}
    \end{equation}
    donde \(n\) es un parámetro que típicamente toma un valor de 3/2 para materiales ferromagnéticos, y \(T_c\) es la temperatura crítica. 